%% sections/hungyu.tex: Hungyu's slides (DQN explorer and the NSGA-II / MAP-Elites baselines). OWNER: Hungyu. Only Hungyu edits this file.
%% - Figures go in figures/hungyu/ and are included as \includegraphics[width=...]{hungyu/<file>}.
%% - Need a macro of your own? Define it here with a hungyu prefix (\newcommand{\hungyuFoo}{...}), never in preamble.tex.
%% - Build just your part: latexmk -pdf solo_hungyu.tex   (or compile solo_hungyu.tex in Overleaf).
%% - Backup slides: put them at the end of this file as \begin{frame}[noframenumbering]{Backup: ...}.
%% - The section's title and cover slide are set in main.tex; ask Renzhi to rename it.


\begin{frame}{DQN, NSGA-II and MAP-Elites: method}
  \textbf{Hungyu Lin, Group 7}
  \begin{itemize}
    \item With 50 to 600 simulations, how well can we map the circuit designs that meet a given spec?
    \item I trained a Double DQN to score candidate designs using the map's predictions and uncertainty. It selects one design per step. During training, a drop in prediction error near the pass/fail boundary earns a positive reward.
    \item I implemented and evaluated two comparison methods. NSGA-II searches for trade-offs across specs. MAP-Elites groups designs by spec margins and keeps the design with the largest minimum margin in each bin.
  \end{itemize}
\end{frame}

\begin{frame}{DQN, NSGA-II and MAP-Elites: training setup}
  \begin{columns}[T]
    \begin{column}{0.5\textwidth}
      \begin{itemize}
        \item I tried six settings with training seed 0 and 250,000 steps per run.
        \item In a 6,200-step pilot sample, about half the rewards had magnitude below $10^{-4}$. I scaled the reward inside DQN.
        \item I kept the $1{,}000\times$ scale and $\gamma=0.9$, then trained three models for one million steps each with different seeds.
      \end{itemize}
    \end{column}
    \begin{column}{0.5\textwidth}
      \centering
      \small
      \textbf{Validation prediction error}\\
      Difference from the \rulename\\[0.6em]
      \begin{tabular}{@{}lr@{}}
        \toprule
        DQN setting & Difference \\
        \midrule
        Default & $+0.265$ \\
        Reward $\times 1{,}000$ & $+0.018$ \\
        Same scale, $\gamma=0.9$ & $+0.002$ \\
        \bottomrule
      \end{tabular}\\[0.4em]
      Three of six settings shown.\\
      Lower error is better.
    \end{column}
  \end{columns}
\end{frame}

\begin{frame}{DQN, NSGA-II and MAP-Elites: results}
  \centering
  \begin{tabular}{lcc}
    \toprule
    Method & Pass/fail (edge-tolerant IoU) $\uparrow$ & \prederr{} $\downarrow$ \\
    \midrule
    Uncertainty rule (reference) & $0$ & $0$ \\
    DQN (three seeds pooled) & $-0.040$ & $+0.013$ \\
    NSGA-II & $-0.217$ & $+0.134$ \\
    MAP-Elites & $-0.324$ & $+0.164$ \\
    \bottomrule
  \end{tabular}

  \vspace{1em}
  Each entry is the method's score minus the rule's, averaged over 19 achievable validation tasks and five simulation budgets.\\
  DQN beat both evolutionary methods on both measures. The reference rule picks the design whose estimated chance of passing is closest to 50\%.
\end{frame}

\begin{frame}{DQN, NSGA-II and MAP-Elites: takeaways}
  \begin{itemize}
    \item Reward scaling alone closed most of the tuning run's prediction-error gap. On stored validation data, all three final DQN models still had higher error than the rule (by $0.002$ to $0.030$).
    \item On six achievable specs for one 15-knob circuit, fresh SPICE runs put DQN's prediction error $0.009$ above the rule's and its pass/fail score $0.042$ below.
    \item I tested both evolutionary methods on SKY130 and GF180. On both processes, their average prediction error exceeded the rule's across the scored tasks. DQN testing is scheduled for Nov.\ 30--Dec.\ 3, after the Nov.\ 29 freeze.
  \end{itemize}
\end{frame}
